\documentclass[11pt,a4paper]{article}
\usepackage{xeCJK}
\usepackage{fontspec}
\usepackage{xcolor}
\usepackage{fancyhdr}
\usepackage{booktabs}
\usepackage{array}
\usepackage{amssymb}
\usepackage{geometry}
\usepackage[protrusion=true]{microtype}
\geometry{margin=1.8cm,top=2cm,headheight=15pt}
\linespread{1.05}

\setmainfont{Baskerville}
\setCJKmainfont{Songti SC}
\newfontfamily{\starfont}{Songti SC}

\definecolor{wrongred}{RGB}{198,40,40}
\definecolor{wrongbg}{RGB}{253,235,233}
\definecolor{rulegray}{RGB}{200,200,200}

\setlength{\emergencystretch}{3em}
\setlength{\fboxsep}{3pt}
\setlength{\parindent}{0pt}

\newcommand{\blank}{\underline{\hspace{2.8cm}}}
\newcommand{\checkmarkbox}{\fbox{\rule{0pt}{0.8ex}\rule{0.8ex}{0pt}}}
\newcommand{\STAR}{\textcolor{wrongred}{\starfont ★}}
% 原卷作答（红色）
\newcommand{\W}[1]{\textcolor{wrongred}{#1}}
% 标准答案（加粗）
\newcommand{\A}[1]{\textbf{#1}}
% 判定标记
\newcommand{\OK}{\textcolor{wrongred}{\ensuremath{\checkmark}}}
\newcommand{\NO}{\textcolor{wrongred}{$\times$}}
% 表格正文列：左对齐，避免 p 列两端对齐造成的 Underfull
\newcolumntype{L}[1]{>{\raggedright\arraybackslash}p{#1}}

\pagestyle{fancy}
\fancyhf{}
\fancyhead[L]{\footnotesize 英语短语默写 小蓝 P31--40（错题订正）}
\fancyhead[R]{\footnotesize 赵文瀚}
\fancyfoot[C]{\thepage}
\renewcommand{\headrulewidth}{0.4pt}
\renewcommand{\footrulewidth}{0pt}

\begin{document}

\begin{center}
{\LARGE\bfseries 英语短语默写 小蓝 P31--40}\\[0.25em]
{\large 错题订正与复盘}\\[0.35em]
{\footnotesize 2029 届高一英语 · 小蓝短语默写 P31--40（共 25 条）· 2026-10}
\end{center}

\vspace{-0.2em}
\noindent{\color{rulegray}\rule{\textwidth}{0.8pt}}

\vspace{0.5em}
\noindent\fcolorbox{black!15}{black!4}{%
  \parbox{\dimexpr\linewidth-2\fboxsep-2\fboxrule\relax}{%
    \small 日期：2026-10\qquad 来源：小蓝短语默写 P31--40（共 25 条）\qquad 姓名：赵文瀚\\[0.3em]
    判错：\textcolor{wrongred}{\textbf{9}}\,/\enspace 25\qquad
    判对：\textbf{16}\,/\enspace 25\qquad
    重做日期：\underline{\hspace{2.4cm}}}}

\vspace{0.6em}
\noindent{\footnotesize 说明：下表按原卷题号列全 25 条。\W{红色}＝原卷作答，\A{加粗}＝标准答案；\OK\ 表示按原卷红笔判对，\NO\ 表示按原卷红笔判错。\STAR\ 表示需优先复习的错项。}

\vspace{0.4em}
\renewcommand{\arraystretch}{0.96}
\noindent{\footnotesize\begin{tabular}{@{}L{0.72cm}@{\hspace{0.08cm}}L{2.03cm}@{\hspace{0.08cm}}L{1.45cm}@{\hspace{0.08cm}}L{4.6cm}@{\hspace{0.08cm}}L{4.7cm}@{\hspace{0.08cm}}L{0.75cm}@{\hspace{0.08cm}}L{1.55cm}@{\hspace{0.08cm}}L{0.9cm}@{}}
\toprule
\textbf{题号} & \textbf{中文短语} & \textbf{提示词} & \textbf{我的作答} & \textbf{标准答案} & \textbf{判定} & \textbf{错误类型} & \textbf{二刷} \\
\midrule
1  & 假笑 & — & an artificial smile & — & \OK & — & \checkmarkbox \\
2  & 尽管尽了最大努力仍不能搬动石头 \STAR & as & \W{As I tried my best, I couldn't move that rock} & \A{As hard as I tried, I couldn't move the rock} & \NO & 句式错误 & \checkmarkbox \\
3  & 他说葡萄牙语好像一个葡萄牙人 \STAR & as & \W{He says Portuguese as if ...（涂改）} & \A{He speaks Portuguese as if he were Portuguese.} & \NO & 用词＋句式 & \checkmarkbox \\
4  & 存点钱 & aside & set aside some money & — & \OK & — & \checkmarkbox \\
5  & 熟睡 \STAR & sound & \W{sound awake} & \A{be sound asleep} & \NO & 反义词误用 & \checkmarkbox \\
6  & 布置某人某事／布置某人做某事 & assign & assign sb. sth.; assign sb. to do sth. & — & \OK & — & \checkmarkbox \\
7  & 不想和这个投资项目有任何瓜葛 \STAR & associate & \W{don't want to associate with the investment program} & \A{don't want to be associated with the investment project} & \NO & 语态／搭配 & \checkmarkbox \\
8  & 在清晨／中午／黄昏／半夜 & — & at dawn / noon / dusk / midnight & — & \OK & — & \checkmarkbox \\
9  & 把经济发展放在第一位 \STAR & attach & \W{attach economic development in the first place} & \A{attach importance to economic development} & \NO & 固定搭配 & \checkmarkbox \\
10 & 详情备查 & — & Further information is available on request & — & \OK & — & \checkmarkbox \\
11 & 自动驾驶汽车 & — & autonomous car / vehicle & — & \OK & — & \checkmarkbox \\
12 & 试图做某事 & attempt & make an attempt to do sth. & — & \OK & — & \checkmarkbox \\
13 & 很多观众 \STAR & — & \W{many audience} & \A{a large audience} & \NO & 名词搭配 & \checkmarkbox \\
14 & 外科大夫今天下午有空吗？ & — & Is the doctor available this afternoon? & — & \OK & — & \checkmarkbox \\
15 & 完全没有睡意 & — & wide awake & — & \OK & — & \checkmarkbox \\
16 & 授予奖品给某人 & award & award sth. to sb. & — & \OK & — & \checkmarkbox \\
17 & 人人都该知道其中风险 \STAR & aware & \W{Everyone should be aware of the endurance} & \A{Everyone should be aware of the risks} & \NO & 用词错误 & \checkmarkbox \\
18 & 感到惭愧 & — & feel ashamed & — & \OK & — & \checkmarkbox \\
19 & 理学学士 & — & Bachelor of Science & — & \OK & — & \checkmarkbox \\
20 & 把……烧为灰烬 & — & burn sth. into ashes & — & \OK & — & \checkmarkbox \\
21 & 学士学位 \STAR & bachelor & \W{degree of bachelor / bachelor degree} & \A{a bachelor's degree} & \NO & 固定搭配 & \checkmarkbox \\
22 & 背弃，抛弃 \STAR & back & \W{put back …; turn one's back …（涂改）} & \A{turn one's back on sb./sth.} & \NO & 固定搭配 & \checkmarkbox \\
23 & 平均来说 & — & on average & — & \OK & — & \checkmarkbox \\
24 & 对……有权威 & — & have authority over sb. & — & \OK & — & \checkmarkbox \\
25 & 向某人保证 & — & assure sb. of sth. & — & \OK & — & \checkmarkbox \\
\bottomrule
\end{tabular}}

\vspace{0.25em}
\noindent{\footnotesize\textbf{原卷旁注}：第 3 题句中多处涂改，完整原答无法确证，表中只录可辨片段；第 4 题有两个尝试，采用可辨的 \emph{set aside some money}；第 5 题原答可辨为 \emph{sound awake}，红笔改为 \emph{sound asleep}；第 2、3、5、7、9、13、17、21、22 题按红笔判错。第 20 题原卷动词笔迹不清，表中按可辨词形暂录；第 22 题除旁注所示外，原答亦有涂改。红笔痕迹与文字重叠处以存疑旁注保留。}

\newpage
\begin{center}
{\Large\bfseries 错因归类与复盘}
\end{center}
\vspace{-0.4em}
\noindent{\color{rulegray}\rule{\textwidth}{0.8pt}}

\vspace{0.5em}
\noindent 说明：按错误性质归类 9 条错题，供汇总对账；逐条二刷结果见第 1 页「二刷」列。

\vspace{0.5em}
\renewcommand{\arraystretch}{1.35}
\begin{center}
\begin{tabular}{@{}p{4.6cm} p{10.3cm} c@{}}
\toprule
\textbf{错误类型} & \textbf{题号} & \textbf{条数} \\
\midrule
句式／固定搭配错误 & {2, 3, 7, 9, 21, 22} & 6 \\
词汇与名词搭配错误 & {5, 13, 17} & 3 \\
\midrule
\textbf{合计} & \textbf{9 条} & \textbf{9} \\
\bottomrule
\end{tabular}
\end{center}

\vspace{0.5em}
\noindent\textbf{关联知识点标签}
\vspace{0.25em}

\noindent\begin{tabular}{@{}p{8.2cm}@{\hspace{0.4cm}}p{8.2cm}@{}}
\begin{minipage}[t]{\linewidth}
\small
\STAR\ \textbf{as 句式}：\textbf{as hard as} + 从句；\textbf{as if} 引导方式状语从句（虚拟语气按语境使用 were）\\[0.2em]
\STAR\ \textbf{sound asleep}：熟睡；不要把 \textbf{awake}（醒着）混入\\[0.2em]
\STAR\ \textbf{associate}：be associated \textbf{with} sth.（与……有关联）\\[0.2em]
\STAR\ \textbf{attach importance to}：重视；importance 后接介词 \textbf{to}
\end{minipage} &
\begin{minipage}[t]{\linewidth}
\small
\STAR\ \textbf{audience}：集合名词，可说 \textbf{a large audience}，不说 many audience\\[0.2em]
\STAR\ \textbf{aware of}：be aware of the \textbf{risks}（意识到风险）\\[0.2em]
\STAR\ \textbf{学位与 back 短语}：\textbf{a bachelor's degree}；\textbf{turn one's back on} sb./sth.\\[0.2em]
\STAR\ \textbf{范围}：小蓝短语 P31--40（本卷 25 条，按原卷题号全列）
\end{minipage} \\
\end{tabular}

\vspace{0.65em}
\noindent\fcolorbox{black!15}{black!4}{%
  \parbox{\dimexpr\linewidth-2\fboxsep-2\fboxrule\relax}{%
    \textbf{复盘记录}\\[0.35em]
    最薄弱的板块：\underline{\hspace{9.0cm}}\\[0.55em]
    主要错误原因：\checkmarkbox 固定搭配不牢\quad \checkmarkbox 介词记混\quad \checkmarkbox 漏词/多词\quad \checkmarkbox 用词/拼写\quad \checkmarkbox 空白未答\\[0.55em]
    本次复习的改进措施：\underline{\hspace{10.1cm}}\\[0.55em]
    \hspace{3.2em}\underline{\hspace{10.1cm}}\\[0.55em]
    下次复习日期：\underline{\hspace{2.2cm}}\qquad
    当前掌握度：\checkmarkbox 仍薄弱\quad \checkmarkbox 基本掌握\quad \checkmarkbox 已掌握
  }}

\end{document}
